% Generated draft: requires mathematical and source review.
\documentclass[11pt]{article}
\usepackage[margin=25mm]{geometry}
\usepackage{fontspec}
\setmainfont{Latin Modern Roman}
\usepackage{amsmath,amssymb,mathtools,unicode-math}
\setmathfont{Latin Modern Math}
\usepackage{enumitem,needspace,xurl,hyperref}
\setlength{\parindent}{0pt}
\setlength{\parskip}{5pt}
\setlength{\emergencystretch}{3em}
\newcommand{\N}{\mathbb{N}}
\newcommand{\Z}{\mathbb{Z}}
\newcommand{\Q}{\mathbb{Q}}
\newcommand{\R}{\mathbb{R}}
\newcommand{\C}{\mathbb{C}}
\newcommand{\F}{\mathbb{F}}
\newcommand{\A}{\mathbb{A}}
\newcommand{\PP}{\mathbb{P}}
\newcommand{\E}{\mathbb{E}}
\newcommand{\eps}{\varepsilon}
\newcommand{\dd}{\,\mathrm{d}}
\begin{document}
\section*{UnsolvedMath ID: 20002351}
\textbf{Source problem code:} catalogue entry
\subsection{Definitions and mathematical statement}
Question 14 asks for a Diophantine injection $\Q^2\hookrightarrow\Q$, motivated by coding blocks of universal quantifiers. The supplied research assessment identifies the intended full-domain algebraic target as
\[
 \exists P\in\Q[X,Y]\ \forall a,b,c,d\in\Q\quad
 \bigl(P(a,b)=P(c,d)\ \Longrightarrow\ a=c\ \wedge\ b=d\bigr).
\]
A polynomial map has a Diophantine graph, but a function with an arbitrary positive-existentially definable graph is a more general notion of ``Diophantine injection.'' The original wording does not expressly choose between these notions. Neither means merely an injection on a restricted subset of $\Q^2$. The aid reports that the polynomial question remains open and that every polynomial of total degree at most two fails to be injective on $\Q^2$; thus a polynomial solution would have degree at least three.

The source credits Jeroen Demeyer with the observation that a polynomial bijection $\N^2\to\N$ permits replacement of each consecutive universal block in a positive prenex formula over $\N$ by one universal quantifier, at the cost of additional existential quantifiers. The cited natural-number coding is from Davis, \emph{Math. Monthly} 80, 236--237. Its essential logical feature is a total surjective decoding map, not just injectivity of a coding map.

Precisely, let $M$ be a nonempty structure and let $r\ge1$. A positive-existential formula is built from atomic formulas by conjunction, disjunction and existential quantification. Suppose $D(y,\bar x)$ is such a formula defining a function $M\to M^r$, so
\[
 (\forall y\in M)(\exists!\bar x\in M^r)D(y,\bar x),
 \qquad
 (\forall\bar x\in M^r)(\exists y\in M)D(y,\bar x).
\]
For any formula $\varphi(\bar x,\bar a)$ one then has
\[
 \forall\bar x\,\varphi(\bar x,\bar a)
 \quad\Longleftrightarrow\quad
 \forall y\,\exists\bar x\,
       \bigl(D(y,\bar x)\wedge\varphi(\bar x,\bar a)\bigr).
\]
Expanding the existential witnesses in $D$ preserves the positive-prenex form and replaces the $r$ universal variables of that block by $y$.

One sufficient way to obtain this decoder is a total injective map $e:M^r\to M$ with positive-existential graph $G(\bar x,y)$, a positive-existential definition $C(y)$ of $M\setminus e(M^r)$, and a named or positive-existentially specified default tuple $\bar c$. Then
\[
 D(y,\bar x)\ \Longleftrightarrow\
 G(\bar x,y)\ \vee\ \bigl(C(y)\wedge\bar x=\bar c\bigr)
\]
is a total surjective decoder. Without the noncode clause, the direct inverse-code rewrite already fails for the tautology $\forall\bar x\,(\bar x=\bar x)$ when $e$ is not onto: its proposed replacement would falsely require every $y$ to lie in the image. This is a limitation of that direct construction, not a proof that no other coding reduction can follow from a bare injection.

The source additionally says that observations in \emph{C. R. Acad. Sci. Paris} 328, 3--8 (1999) imply the desired injection from a generalized $abc$ conjecture. No exact conjecture, theorem or bibliographic identification beyond that citation is supplied, so this remains a reported conditional connection rather than an instantiated theorem.
\subsection{Short English statement}
Seek a polynomial injection from all rational pairs into the rationals. The research aid leaves that question open and gives a precise sufficient decoder condition for collapsing universal quantifier blocks; a non-surjective injection alone does not justify the direct inverse-code rewrite.
\subsection{Sources}
\begin{itemize}
\item Cornelissen, Question 14, American Institute of Mathematics workshop \emph{Extensions of Hilbert's Tenth Problem}, \url{https://aimath.org/WWN/hilberts10th/hilberts10th.pdf}. The coding observation is attributed there to Jeroen Demeyer.
\item Davis, \emph{Math. Monthly} 80, 236--237, as cited in the question for a natural-number polynomial bijection; \emph{C. R. Acad. Sci. Paris} 328, 3--8 (1999), as cited for the generalized-$abc$ connection. Titles and the author of the latter reference are not supplied.
\item Local UnsolvedMath registry, permanent ID 20002351, source code \texttt{catalogue entry}, canonical location \texttt{aim-logic-notes.json}, \texttt{notes[126]}, record dated 14 August 2026. The original literature field is empty.
\item Supplied research attempt 1: partial result, problem status at run open, rigor labelled proved, novelty confidence low. The decoder criterion and degree bound are machine-generated research aid, not peer reviewed, requiring independent verification before citation.
\end{itemize}
\end{document}
